\crefalias{section}{appendix}
\section{The Geometric Origin of Quantum Phenomena: Substrate Mechanics at the Hardware Limit}
\label{app:quantum_fundamentals}

By applying the formalism of Entropic Action ($S_\Phi = \int \mathcal{D} d\mu$) we can re-derive superposition, wave-function collapse, and non-local entanglement as the statistical mechanics of substrate survival and projection artifacts.

\subsection{The Survivorship Filter and Quantum Mechanics}

In standard Quantum Field Theory, the Feynman Path Integral is postulated, relying on Wick Rotation to transform oscillatory quantum phase into a decaying thermodynamic partition function. In the $E_8$-Persistence theory, this hierarchy is inverted: the discrete, thermodynamic substrate is the exact physical reality, and the continuous complex path integral is the macroscopic approximation (as established in \cref{sec:entropic_lagrangian_density}).

Persistence is governed by the \textbf{Survivorship Filter}. The probability of a structural history $\Gamma$ maintaining causal coherence against the substrate's noise floor is:
\begin{equation}
    P(\text{survival} \mid \Gamma) \propto e^{-\beta S_\Phi[\Gamma]}
\end{equation}
where $\beta = 1/C_{\text{channel}}$ represents the inverse entropic tolerance of the substrate, a dimensionless geometric invariant fixed by the lattice's self-dual resonance.

This survivorship filter naturally recovers the Born rule. For a state $|\psi\rangle = \sum_k c_k |k\rangle$ with geometric support over multiple internal addresses, the entropic action of trajectory $\Gamma_k$ terminating at address $k$ is:
\begin{equation}
    S_\Phi[\Gamma_k] = -\ln |c_k|^2 + S_{\text{base}}
\end{equation}
where $S_{\text{base}}$ is path-independent. The survival probability is therefore:
\begin{equation}
    P(k) \propto e^{-\beta S_\Phi[\Gamma_k]} = |c_k|^{2\beta} e^{-\beta S_{\text{base}}}
\end{equation}
With $\beta = 1$ (the geometric invariant fixing the survivorship filter to the lattice's self-dual resonance), normalization yields $P(k) = |c_k|^2$. The Born rule is thus derived as the thermodynamic equilibrium of the Survivorship Filter.

\subsection{Entanglement as Unprojected Adjacency}

Quantum entanglement represents a non-local correlation across macroscopic spacetime distances, seemingly violating special relativity. Standard physics treats this as a purely algebraic consequence of the Hilbert space tensor product. The $E_8$-Persistence theory resolves this paradox geometrically by modeling entanglement as \textbf{Unprojected Adjacency} within the 8-dimensional substrate.

The $E_8 \to D_4 \oplus D_4$ projection strictly partitions the lattice into an external spacetime manifold ($x^\mu \in \text{Sector A}$) and an orthogonal internal symmetry space ($\theta^a \in \text{Sector B}$). For two entangled particles sharing an internal state, their geometric distance on the full substrate is the sum of their orthogonal metric distances:
\begin{equation}
    ds^2_{\text{total}} = \eta_{\mu\nu}\Delta x^\mu \Delta x^\nu + \delta_{ab}\Delta \theta^a \Delta \theta^b
\end{equation}
While the particles may be kinematically distant in the observable $D=4$ spacetime manifold ($\Delta x^\mu \gg 0$), their shared internal configuration dictates they remain identically colocated in the internal sector ($\Delta \theta^a = 0$). 

\textbf{Preservation of Causality (No FTL Signaling):} Because information (energy transfer) propagates at the bandwidth limit $c_{\text{eff}} = 1$ strictly \textit{within the projected causal manifold} (Sector A), special relativity is perfectly preserved in spacetime. Alice cannot send a faster-than-light macroscopic message to Bob.

\textbf{Resolution of Non-Locality (Bell's Inequality):} Crucially, as established in Paper I, the metric projection assigns the active temporal generator ($\partial_t$) exclusively to the spacetime manifold. The internal symmetry space (Sector B) retains a strictly Euclidean signature. Because it possesses no temporal generator, it is perfectly static relative to the spacetime metric. 

When a measurement on Alice's particle forces an Address Resolution (see Section \ref{subsec:measurement}), the required geometric update to the shared internal state does not ``travel'' through time; the collapse is geometrically instantaneous. From the perspective of the 4D observer, the correlation appears non-local. From the perspective of the $E_8$ substrate, the two particles never separated, and their shared internal update requires zero clock cycles to resolve.

\subsection{Measurement and Basis Selection}
\label{subsec:measurement}

Wave-particle duality and measurement collapse are direct consequences of the \textbf{Equipartition of Entropic Cost} and the node's finite channel capacity ($\nu = 16$).

When a particle passes through a double-slit apparatus, it occupies a single subspace address with geometric support over multiple trajectories in the internal $D_4$ sector. To minimize its transition flux ($f$) and avoid the dissipation penalty of the Survivorship Filter, the particle propagates its structural complexity ($\mathbf{K}$) as a delocalized phase front. 

\textbf{Address Resolution (Capacity Exhaustion):} Measurement is an interaction that demands Address Resolution. Placing a detector at the slits forces the physical lattice to evaluate the particle's exact spatial coordinate. Resolving this localized fermion node intrinsically consumes exactly 3 units of the dimensional budget, saturating the local node's physical bit-depth ($\nu = 16$). 

The local node physically lacks the capacity to simultaneously encode the high-resolution spatial address (the ``Which-Path'' measurement) and maintain the coherent internal phase array (the geometric support over the unprojected $D_4$ internal symmetry sector). To avoid information overflow, the substrate must sacrifice the phase coherence to maintain unitary solvency at the node. This forces a discrete projection basis, irreversibly destroying the phase coherence required for interference and generating localized Entropic Action ($\Delta S_\Phi \ge \ln 2$).

Formally, this process implements a projection-valued measure $\{P_k\}$ on the spacetime factor: $P_k = |x_k\rangle\langle x_k| \otimes \mathbb{I}_{\text{internal}}$. The post-measurement state is $|\psi'\rangle = P_k |\psi\rangle / \sqrt{\langle\psi|P_k|\psi\rangle}$, with the internal residual $|\phi_k\rangle_{\text{internal}}$ decoupled from spacetime propagation by the Euclidean signature of Sector B.

\textbf{Delayed-Choice and Basis Selection:} The ``delayed choice'' simply dictates \textit{which projection basis} the macroscopic apparatus forces the substrate to use at the final evaluation step. If distinguishing information is rigidly preserved in the environment, the local dissipation density $\mathcal{D}$ forces projection onto the localized spacetime eigenstates (Slit A vs. Slit B). If the information is ``erased,'' the energetic demand for exact address resolution is lifted, and the substrate defaults to the minimal-action superposition basis (A+B). The erasure is purely a physical reorientation of the local projection operator occurring strictly in forward time.